\chapter{Growing Pains}

\section*{Definition}
Growing pains are recurrent, bilateral leg pain in children aged 3--12 years, occurring in the late afternoon or evening and resolving by morning. It is a diagnosis of exclusion with no identifiable organic pathology, and episodes typically cease by adolescence.

\section*{What Happens in Your Body}
The exact mechanism is unknown, though several theories exist. One hypothesis suggests that rapid bone growth outpaces the elongation capacity of periarticular soft tissues (muscles, tendons, fascia), creating relative tension and discomfort. Another implicates altered pain threshold or central pain processing, as affected children often have lower pain tolerance. There is no evidence of inflammatory, infectious, or neoplastic processes.

\section*{Causes}
\begin{itemize}
  \item \textbf{Idiopathic:} No known cause in the majority of cases
  \item \textbf{Musculoskeletal fatigue:} Increased physical activity, poor postural habits, hypermobility
  \item \textbf{Genetic:} Family history of growing pains or restless legs syndrome is common
  \item \textbf{Biomechanical:} Flat feet, genu valgum, or tibial torsion may contribute
  \item \textbf{Psychosocial:} Stress, anxiety, or secondary gain may influence symptom reporting
\end{itemize}

\section*{When to See a Doctor}
See your doctor if:
\begin{itemize}
  \item Pain is unilateral, persistent during the day, or wakes the child from sleep
  \item Associated with fever, limping, refusal to bear weight, or joint swelling
  \item Pain localized to a joint rather than the muscle belly
  \item Systemic symptoms: fatigue, rash, weight loss, or pallor
\end{itemize}

\section*{Related Symptoms}
\begin{itemize}
  \item Bilateral aching of the thighs, calves, and popliteal fossa
  \item Onset in late afternoon or evening after active play
  \item Complete resolution by morning with no residual pain or stiffness
  \item Normal activity and gait during the day
  \item No associated joint swelling, erythema, or fever
\end{itemize}
\textbf{Important:} Growing pains are a benign self-limited condition; red flags include persistent morning pain, joint swelling, or systemic symptoms that warrant investigation for juvenile idiopathic arthritis or malignancy.

\section*{Self-Care / Home Management}
\begin{itemize}
  \item Reassurance that pain is benign and self-limited
  \item Gentle massage of the affected muscles
  \item Warm compresses or a warm bath before bedtime
  \item Stretching exercises for the quadriceps, hamstrings, and calf muscles during the day
  \item Acetaminophen or ibuprofen as needed for moderate-to-severe episodes
\end{itemize}

\section*{How Your Doctor Will Evaluate You}
\begin{itemize}
  \item Clinical diagnosis based on well-established criteria: bilateral leg pain, evening onset, no functional limitation, normal examination
  \item Normal gait, full range of motion, and absence of joint swelling or tenderness on examination
  \item No laboratory testing or imaging needed in typical cases
  \item Consider ESR, CRP, CBC, and radiography if red flags are present
\end{itemize}

\section*{Treatment Options}
\begin{itemize}
  \item Parent and child education with reassurance about benign nature
  \item Analgesics (acetaminophen or ibuprofen) only during symptomatic episodes
  \item Evening stretching program targeting the lower extremity muscle groups
  \item Magnesium or calcium supplementation (limited evidence of benefit)
  \item No restriction of daytime physical activity
  \item Follow-up if symptoms persist beyond adolescence or change character
\end{itemize}

\clearpage
